
\subsection{$A_K$ 同态的构造 I}
\label{subsec:7-5-4}

在第~\ref{sec:7-5} 节的余下部分，
我们只讨论 $A_\infty$ 代数，而不讨论滤过 $A_\infty$ 代数。
因此，下文以 $\mathfrak m_k$、$\mathfrak f_k$
代替 $\overline{\mathfrak m}_k$、$\overline{\mathfrak f}_k$。

为证明定理~\ref{thm:7-5-1}，
我们将构造一个从第~\ref{subsec:7-2-5} 小节所得 $A_K$ 代数
到德拉姆复形 $\Omega(L)$ 的 $A_K$ 同态
\[
\mathfrak f=\{\mathfrak f_i\}_{i=1}^{K}.
\]
先构造 $\mathfrak f_1$。
它就是通常的\emph{磨光子光滑化}（见 \textbf{[dRh60]}）。
具体地，在 $L\times L$ 上取一个闭 $n$-形式
$\mathfrak w_\epsilon$；它支撑在对角线
$L\subset L\times L$ 的一个 $\epsilon$-邻域内，
并且其德拉姆上同调类是该对角线的庞加莱对偶。
定义
\begin{equation}
\mathfrak f_1(P)
=\left.\pi_{2!}\right|_{P\times L}
\left(\left.\mathfrak w_\epsilon\right|_{P\times L}\right).
\tag{7.5.13}\label{eq:7-5-13}
\end{equation}
这里 $\pi_2$ 是到第二个因子的投影，
$\pi_{2!}$ 表示沿纤维积分。

\begin{lemma}{7.5.14}\label{lem:7-5-14}
式~\eqref{eq:7-5-13} 给出的映射 $\mathfrak f_1$
定义了从 $(\mathbb R\mathcal X_L,\delta)$ 到 $\Omega(L)$
的上链同态，并在上同调上诱导同构。
\end{lemma}

引理~\ref{lem:7-5-14} 的证明是熟知的
（见 \textbf{[dRh60]}）。

下面在定义~\ref{def:7-5-6} 的意义下，
给出 $\mathfrak f_1(P)$ 的一个德拉姆表示。
置
\[
\mathcal N'_2=\{\mathrm{point}\},
\qquad
U_2(\mathfrak f)=U_2\times\mathcal N'_2.
\]
这里采用第~\ref{subsec:7-5-3} 小节末尾所述的记号；
一般 $\mathcal N'_{k+1}$ 的定义见后面的命题 7.5.43。
把 $U_2$ 上的丛 $E_2$ 沿投影
$U_2(\mathfrak f)\to U_2$ 拉回，所得丛记作
$E_2(\mathfrak f)$。

取一个无边界流形 $W_1$，并考虑由 $w\in W_1$
参数化的 $E_2(\mathfrak f)$ 截面族
\[
\mathfrak t_{\mathfrak f,2}
=\{\mathfrak t_{\mathfrak f,2,w}\mid w\in W_1\}.
\]
换言之，$\mathfrak t_{\mathfrak f,2}$ 是丛
\[
W_1\times E_2(\mathfrak f)
\longrightarrow W_1\times U_2(\mathfrak f)
\]
的一个截面。
考察
$\mathfrak t_{\mathfrak f,2}^{-1}(0)
\subset W_1\times U_2(\mathfrak f)$，
并将在 $W_1$ 上取一个光滑概率测度
（见定义~\ref{def:7-5-6}）。
从这里起，$\mathfrak t_{\mathfrak f,k+1}$
表示用来定义 $\mathfrak f_k$ 的截面族，
它是原仓西映射 $s_{\mathfrak f,k+1}$ 的扰动。

置
\[
U_2(\mathfrak f,P)
=U_2(\mathfrak f){}_{ev_1}\!\mathbin{\times_L}P.
\]
作为集合，它同
\[
ev_1^{-1}(P)\cap U_2(\mathfrak f)
\subset U_2(\mathfrak f)
\]
一一对应。
此后，为明确所用的 $ev_1$，
我们把 $P$ 在 $ev_1\colon U_2(\mathfrak f)\to L$
下的逆像写成 $ev_1^{-1}(P)\cap U_2(\mathfrak f)$。
还要指出同构
\[
U_2(\mathfrak f,P)
\cong U_2(\mathfrak f)
\mathbin{\times_{L^2}}(P\times L),
\]
后者作为集合可同
$\mathcal N'_2\times\bigl(U_2\cap(P\times L)\bigr)$ 识别
\translatornote{源书在这个以 $L^2$ 为底的纤维积左侧标作
$ev_1$；但 $ev_1$ 的靶是 $L$，不能作到 $L^2$ 的结构映射。
这里按完整评价映射 $(ev_1,ev_0)$ 所给的纤维积书写。
见疑误 A-0617。}。
以后不再重复这类说明。

\begin{lemma}{7.5.15}\label{lem:7-5-15}
可以选择上述参数空间 $W_1$ 和截面族
$\mathfrak t_{\mathfrak f,2}$，使得
\[
\mathfrak f_1(P)
=ev_{0!}\bigl(\pi_{W_1}^*(\omega_1)\bigr).
\]
这里
$ev_0\colon W_1\times U_2(\mathfrak f,P)\to L$
是 $ev_0\colon W_1\times U_2(\mathfrak f)\to L$ 的限制，
而
$\pi_{W_1}\colon W_1\times U_2(\mathfrak f,P)\to W_1$
是投影。
\end{lemma}

证明很容易，故从略。

现在开始构造 $\mathfrak f_2$。
$(\mathfrak f_1,\mathfrak f_2)$ 为 $A_2$ 同态的条件可写成
\begin{align}
&\mathfrak f_1\bigl(\mathfrak m_2(P_1,P_2)\bigr)
+\mathfrak f_2\bigl(\mathfrak m_1(P_1),P_2\bigr)
+(-1)^{\deg P_1}
 \mathfrak f_2\bigl(P_1,\mathfrak m_1(P_2)\bigr)
\notag\\
&\quad={}
(-1)^{\deg P_1(\deg P_2+1)}
\mathfrak f_1(P_1)\wedge\mathfrak f_1(P_2)
+(-1)^{\deg P_1+\deg P_2-1}
d\bigl(\mathfrak f_2(P_1,P_2)\bigr).
\tag{7.5.16}\label{eq:7-5-16}
\end{align}
我们将按
\[
\dim P_1+\dim P_2
=\dim L-\deg P_1+\dim L-\deg P_2
\]
归纳定义 $\mathfrak f_2(P_1,P_2)$。
归纳时假设式~\eqref{eq:7-5-16} 中除最后一项外的各项
都已经定义，并要求该式成立。
随后，考察给出 $\mathfrak f_2(P_1,P_2)$ 德拉姆表示的
带仓西结构空间的各个边界分支，便可导出关系
\eqref{eq:7-5-16}。

置
\[
\mathcal N'_3=[0,1],
\qquad
U_3(\mathfrak f)=U'_3\times\mathcal N'_3.
\]
这里 $U'_3\subset U_3\subset L^3$ 是对角线
$\Delta\subset L^3$ 的一个开邻域，稍后再确定。
$U_3$ 上的丛 $E_3$，亦即 $L^3$ 中对角线 $\Delta$
的法丛的拉回，进一步拉回为 $U_3(\mathfrak f)$ 上的丛
$E_3(\mathfrak f)$。
定义 $E_3(\mathfrak f)$ 的截面（原仓西映射）
$s_{\mathfrak f,3}$，使其零集为
\[
s_{\mathfrak f,3}^{-1}(0)
=\Delta\times\mathcal N'_3
\cong L\times\mathcal N'_3.
\]
置
\[
U_3(\mathfrak f,P_1,P_2)
=U_3(\mathfrak f){}_{(ev_1,ev_2)}\!\mathbin{\times_{L^2}}
(P_1\times P_2).
\]
现在考察它的边界。
在仓西结构的意义下考察
$U_3(\mathfrak f,P_1,P_2)$ 的边界时，
不计 $U'_3$ 的端
\translatornote{源书先把零集的两个 $\mathcal N'_3$
均误写成 $\mathcal N'_2$，又把下式左端写成
$\partial U'_3(\mathfrak f)$；由
$U_3(\mathfrak f)=U'_3\times\mathcal N'_3$ 及
$\mathcal N'_3=[0,1]$ 可知应如译文。见疑误 A-0618。}，故
\[
\partial U_3(\mathfrak f)
=U'_3\times\partial\mathcal N'_3.
\]
因此
\begin{align}
\partial U_3(\mathfrak f,P_1,P_2)
={}&U_3(\mathfrak f,\partial P_1,P_2)
\cup U_3(\mathfrak f,P_1,\partial P_2)
\notag\\
&\cup
\bigl(U'_3\times\partial\mathcal N'_3\bigr)
{}_{(ev_1,ev_2)}\!\mathbin{\times_{L^2}}(P_1\times P_2).
\tag{7.5.17}\label{eq:7-5-17}
\end{align}
式~\eqref{eq:7-5-17} 右端的前两项
描述式~\eqref{eq:7-5-16} 左端的第二、第三项。

再考察式~\eqref{eq:7-5-17} 的第三项。
回忆 $\partial\mathcal N'_3=\{0,1\}$。
对 $i=0,1$，置
\begin{equation}
\partial_iU_3(\mathfrak f,P_1,P_2)
=\bigl(U'_3\times\{i\}\bigr)
{}_{(ev_1,ev_2)}\!\mathbin{\times_{L^2}}(P_1\times P_2).
\tag{7.5.18}\label{eq:7-5-18}
\end{equation}
我们将证明，式~\eqref{eq:7-5-16} 左、右两端的第一项
分别对应于
$\partial_0U_3(\mathfrak f,P_1,P_2)$ 与
$\partial_1U_3(\mathfrak f,P_1,P_2)$。
为此，需要找出 $s_{\mathfrak f,3}$ 的一个扰动族
$\mathfrak t_{\mathfrak f,3}$。
注意，对 $i=0,1$ 都有
\[
s_{\mathfrak f,3}^{-1}(0)
\cap\partial_iU_3(\mathfrak f,P_1,P_2)
\cong P_1\cap P_2.
\]
这是因为
$s_{\mathfrak f,3}^{-1}(0)\cong L\times\mathcal N'_3$，
且 $\Delta\cap(P_1\times P_2)\cong P_1\cap P_2$。

先考察
$\mathfrak f_1\bigl(\mathfrak m_2(P_1,P_2)\bigr)$
的德拉姆表示。
该表示使用空间
$W_1\times U_3(\mathfrak f_1\circ\mathfrak m_2,P_1,P_2)$，
其中
\begin{equation}
U_3(\mathfrak f_1\circ\mathfrak m_2,P_1,P_2)
=U_3(\mathfrak m,P_1,P_2)
{}_{ev_0}\!\mathbin{\times_{ev_1}}U_2(\mathfrak f).
\tag{7.5.19}\label{eq:7-5-19}
\end{equation}
在这个空间上取障碍丛
$E_3(\mathfrak m)\times E_2(\mathfrak f)$ 的拉回。
二元组
\[
\bigl(U_3(\mathfrak f_1\circ\mathfrak m_2,P_1,P_2),
E_3(\mathfrak m)\times E_2(\mathfrak f)\bigr)
\]
定义一个仓西结构，而
$\mathfrak s_{\mathfrak m,3}\times\mathfrak t_{\mathfrak f,2}$
是 $W_1\times
U_3(\mathfrak f_1\circ\mathfrak m_2,P_1,P_2)$ 上的截面。
有投影
\[
\begin{aligned}
\pi_{\mathfrak f_1\circ\mathfrak m_2,W_1}\colon
&\bigl(\mathfrak s_{\mathfrak m,3}
\times\mathfrak t_{\mathfrak f,2}\bigr)^{-1}(0)
\cap\bigl(W_1\times
U_3(\mathfrak f_1\circ\mathfrak m_2,P_1,P_2)\bigr)
\longrightarrow W_1
\end{aligned}
\]
以及评价映射
\[
\begin{aligned}
ev_{\mathfrak f_1\circ\mathfrak m_2,0}\colon
&\bigl(\mathfrak s_{\mathfrak m,3}
\times\mathfrak t_{\mathfrak f,2}\bigr)^{-1}(0)
\cap\bigl(W_1\times
U_3(\mathfrak f_1\circ\mathfrak m_2,P_1,P_2)\bigr)
\longrightarrow L.
\end{aligned}
\]
于是得到
$\mathfrak f_1\bigl(\mathfrak m_2(P_1,P_2)\bigr)$
的德拉姆表示
\begin{equation}
\mathfrak f_1\bigl(\mathfrak m_2(P_1,P_2)\bigr)
=ev_{\mathfrak f_1\circ\mathfrak m_2,0!}
\bigl(\pi_{\mathfrak f_1\circ\mathfrak m_2,W_1}^*\omega_1\bigr).
\tag{7.5.20}\label{eq:7-5-20}
\end{equation}

下面说明这个表示同
$\partial_0U_3(\mathfrak f,P_1,P_2)$ 的关系。
定义
\[
\operatorname{Comp}\colon
U_3{}_{ev_0}\!\mathbin{\times_{ev_1}}U_2
\longrightarrow L^3
\]
为
\[
\operatorname{Comp}
\bigl((x_0,x_1,x_2),(y_0,y_1)\bigr)
=(x_1,x_2,y_0).
\]
这里 $x_0=y_1$，因为
$ev_0(x_0,x_1,x_2)=ev_1(y_0,y_1)$。
可以选择 $U'_3$，使得
\[
U'_3\subset
\operatorname{Comp}
\bigl(U_3{}_{ev_0}\!\mathbin{\times_{ev_1}}U_2\bigr).
\]
$\operatorname{Comp}$ 是浸没，其纤维维数为
$n=\dim L$；这一自由度恰好对应于选择 $x_0=y_1$。
定义 $\operatorname{Comp}$ 的一个截面
\begin{equation}
\begin{aligned}
\varphi\colon U'_3&\longrightarrow
U_3{}_{ev_0}\!\mathbin{\times_{ev_1}}U_2,\\
\varphi(x_1,x_2,y_0)
&=\bigl((x_0,x_1,x_2),(y_0,x_0)\bigr),
\end{aligned}
\tag{7.5.21}\label{eq:7-5-21}
\end{equation}
其中 $x_0$ 是 $x_1,x_2,y_0$ 的黎曼质心
（见 \textbf{[GrKa73]}）。
也就是说，把 $x_0$ 定义为满足
\begin{equation}
\begin{aligned}
&d(x_0,x_1)+d(x_0,x_2)+d(x_0,y_0)\\
&\qquad={}
\inf\{d(x,x_1)+d(x,x_2)+d(x,y_0)\mid x\in L\}
\end{aligned}
\tag{7.5.22}\label{eq:7-5-22}
\end{equation}
的点。
这里在 $L$ 上固定一个黎曼度量，$d$ 是距离函数。
这种 $x_0$ 的存在唯一性来自 $d$ 的局部凸性，
以及距离 $d(x_1,x_2)$、$d(x_1,y_0)$ 很小这一事实
\translatornote{源书式 (7.5.22) 原样使用距离之和，
并据此断言存在唯一；而通常的黎曼质心构造使用距离平方之和。
现有本地证据不足以判定作者是否漏排平方，故译文保留源式，
不作静默修正。见疑误 A-0619。}。

注意
\[
\mathcal M_3^{\mathrm{main}}\times\mathcal N'_2
\cong\{0\}\subset\partial\mathcal N'_3,
\]
因为
$\mathcal M_3^{\mathrm{main}}\cong\mathcal N'_2
\cong\{\text{一点}\}$。
所以式~\eqref{eq:7-5-21} 诱导浸入
\[
\varphi_{\mathfrak f_1\circ\mathfrak m_2,\mathfrak f_2}
\colon
\partial_0U_3(\mathfrak f,P_1,P_2)
\longrightarrow
U_3(\mathfrak f_1\circ\mathfrak m_2,P_1,P_2).
\]
这里回忆
$U_3(\mathfrak f)=U'_3\times\mathcal N'_3$。

\begin{lemma}{7.5.24}\label{lem:7-5-24}
在
$\varphi_{\mathfrak f_1\circ\mathfrak m_2,\mathfrak f_2}$
之上，存在向量丛的单射同态
\[
E_3(\mathfrak f)
\longrightarrow E_3(\mathfrak m)\times E_2(\mathfrak f),
\]
并且在 $\partial_0U_3(\mathfrak f,P_1,P_2)$ 上有丛同构
\begin{equation}
\frac{
\varphi_{\mathfrak f_1\circ\mathfrak m_2,\mathfrak f_2}^*
\bigl(E_3(\mathfrak m)\times E_2(\mathfrak f)\bigr)}
{E_3(\mathfrak f)}
\cong
\frac{
\varphi_{\mathfrak f_1\circ\mathfrak m_2,\mathfrak f_2}^*
T U_3(\mathfrak f_1\circ\mathfrak m_2,P_1,P_2)}
{T\partial_0U_3(\mathfrak f,P_1,P_2)}.
\tag{7.5.25}\label{eq:7-5-25}
\end{equation}
\end{lemma}

这里不讨论同构~\eqref{eq:7-5-25} 如何影响定向；
参见注~\ref{rem:7-5-11}~(1)。

\begin{proof}
定义 $\operatorname{Cent}\colon U_3\to L$，
使 $x_0=\operatorname{Cent}(x_1,x_2,y_0)$
满足式~\eqref{eq:7-5-22}。
对这样的 $x_0$，沿最短测地线作平行移动，得到同构
\[
T_{x_1}L\cong T_{x_2}L\cong T_{y_0}L\cong T_{x_0}L.
\]
它诱导丛同构
\[
(\operatorname{Cent}^*TL)^{\oplus3}\cong TU_3.
\]
由定义容易看出
\[
E_3(\mathfrak f)
\cong\frac{TU_3}{\operatorname{Cent}^*TL},
\]
其中
$\operatorname{Cent}^*TL\to TU_3$
是在上述同构下由 $v\mapsto(v,v,v)$ 定义的。

利用
\[
T_{(x_1,x_2,y_0)}U_3
\cong T_{x_1}L\oplus T_{x_2}L\oplus T_{y_0}L
\]
把切丛分解为
\[
TU_3\cong T^1U_3\oplus T^2U_3\oplus T^3U_3.
\]
于是
\[
\begin{aligned}
\varphi_{\mathfrak f_1\circ\mathfrak m_2,\mathfrak f_2}^*
E_3(\mathfrak m)
&\cong
\frac{\operatorname{Cent}^*TL\oplus T^1U_3\oplus T^2U_3}
{\operatorname{Cent}^*TL}
\cong T^1U_3\oplus T^2U_3,\\
\varphi_{\mathfrak f_1\circ\mathfrak m_2,\mathfrak f_2}^*
E_2(\mathfrak f)
&\cong
\frac{\operatorname{Cent}^*TL\oplus T^3U_3}
{\operatorname{Cent}^*TL}
\cong T^3U_3.
\end{aligned}
\]
由此得到
$E_3(\mathfrak f)\to E_3(\mathfrak m)\times E_2(\mathfrak f)$。
此外，
\[
\frac{
\varphi_{\mathfrak f_1\circ\mathfrak m_2,\mathfrak f_2}^*
\bigl(E_3(\mathfrak m)\times E_2(\mathfrak f)\bigr)}
{E_3(\mathfrak f)}
\cong\operatorname{Cent}^*TL.
\]
另一方面，
\[
\varphi_{\mathfrak f_1\circ\mathfrak m_2,\mathfrak f_2}^*
T U_3(\mathfrak f_1\circ\mathfrak m_2,P_1,P_2)
\cong\operatorname{Cent}^*TL\oplus TU_3.
\]
事实上，在 $(x_1,x_2,y_0)$ 处，右端各分量是
\[
T_{x_0}L\oplus T_{x_1}L\oplus T_{x_2}L\oplus T_{y_0}L.
\]
因此
\[
\frac{
\varphi_{\mathfrak f_1\circ\mathfrak m_2,\mathfrak f_2}^*
T U_3(\mathfrak f_1\circ\mathfrak m_2,P_1,P_2)}
{T\partial_0U_3(\mathfrak f,P_1,P_2)}
\cong\operatorname{Cent}^*TL.
\]
式~\eqref{eq:7-5-25} 得证。
\end{proof}

利用引理~\ref{lem:7-5-24}，
可以把
$\varphi_{\mathfrak f_1\circ\mathfrak m_2,\mathfrak f_2}$
提升为丛的单射同态
\[
\widetilde\varphi_{\mathfrak f_1\circ\mathfrak m_2,\mathfrak f_2}
\colon E_3(\mathfrak f)
\longrightarrow
\varphi_{\mathfrak f_1\circ\mathfrak m_2,\mathfrak f_2}^*
\bigl(E_3(\mathfrak m)\times E_2(\mathfrak f)\bigr),
\]
使它在商丛上诱导同构~\eqref{eq:7-5-25}。
还可以选取
$\widetilde\varphi_{\mathfrak f_1\circ\mathfrak m_2,\mathfrak f_2}$，
使其保持仓西映射
\translatornote{源书此句写成选取不带波浪号的底空间映射
$\varphi_{\mathfrak f_1\circ\mathfrak m_2,\mathfrak f_2}$
来保持仓西映射；能同截面相容并作坐标变换的应是刚提升所得的
丛映射 $\widetilde\varphi$，故译文补回波浪号。
见疑误 A-0621。}。
下面说明怎样用这对映射定义仓西坐标图的坐标变换。

回忆，$\partial_0U_3(\mathfrak f,P_1,P_2)$
的原仓西邻域由丛
\begin{equation}
[0,\rho)\times E_3(\mathfrak f)
\longrightarrow
[0,\rho)\times\partial_0U_3(\mathfrak f,P_1,P_2)
\tag{7.5.26.1}\label{eq:7-5-26-1}
\end{equation}
给出。
取它的另一个仓西坐标图
\begin{equation}
[0,\rho)\times E_3(\mathfrak m)\times E_2(\mathfrak f)
\longrightarrow
[0,\rho)\times
U_3(\mathfrak f_1\circ\mathfrak m_2,P_1,P_2).
\tag{7.5.26.2}\label{eq:7-5-26-2}
\end{equation}
从式~\eqref{eq:7-5-26-1} 到
式~\eqref{eq:7-5-26-2} 的坐标变换由丛映射对
\[
\bigl(
\widetilde\varphi_{\mathfrak f_1\circ\mathfrak m_2,\mathfrak f_2},
\varphi_{\mathfrak f_1\circ\mathfrak m_2,\mathfrak f_2}
\bigr)
\]
给出。

接下来考察式~\eqref{eq:7-5-16} 右端第一项
$\mathfrak f_1(P_1)\wedge\mathfrak f_1(P_2)$ 的德拉姆表示。
算子
$(P_1,P_2)\mapsto
\mathfrak f_1(P_1)\wedge\mathfrak f_1(P_2)$
使用模空间
\[
(W_1\times W_1)\times
U_3(\mathfrak f_1\wedge\mathfrak f_1,P_1,P_2),
\]
其中
\[
U_3(\mathfrak f_1\wedge\mathfrak f_1,P_1,P_2)
:=U_2(\mathfrak f,P_1)
{}_{ev_0}\!\mathbin{\times_{ev_0}}U_2(\mathfrak f,P_2).
\]
按定义有同构
\[
\begin{aligned}
&(W_1\times W_1)\times
U_3(\mathfrak f_1\wedge\mathfrak f_1,P_1,P_2)\\
&\qquad\cong
\bigl(W_1\times U_2(\mathfrak f,P_1)\bigr)
{}_{ev_0}\!\mathbin{\times_{ev_0}}
\bigl(W_1\times U_2(\mathfrak f,P_2)\bigr).
\end{aligned}
\]
利用这个同构，把丛
$E_2(\mathfrak f)\times E_2(\mathfrak f)$
拉回到
$(W_1\times W_1)\times
U_3(\mathfrak f_1\wedge\mathfrak f_1,P_1,P_2)$，
得到带截面
$\mathfrak t_{\mathfrak f,2}\wedge\mathfrak t_{\mathfrak f,2}$
的丛 $E(\mathfrak f_1\wedge\mathfrak f_1)$。
这里，截面
$\mathfrak t_{\mathfrak f,2}\wedge\mathfrak t_{\mathfrak f,2}$
是在上述同构下，把
\[
\bigl(W_1\times U_2(\mathfrak f,P_1)\bigr)
{}_{ev_0}\!\mathbin{\times_{ev_0}}
\bigl(W_1\times U_2(\mathfrak f,P_2)\bigr)
\]
上的截面
$\mathfrak t_{\mathfrak f,2}\times\mathfrak t_{\mathfrak f,2}$
拉回所得。

有投影
\[
\begin{aligned}
\pi_{\mathfrak f_1\wedge\mathfrak f_1,W_1\times W_1}\colon
&\bigl(\mathfrak t_{\mathfrak f,2}
\wedge\mathfrak t_{\mathfrak f,2}\bigr)^{-1}(0)\\
&\quad\cap
\bigl(W_1\times W_1\times
U_3(\mathfrak f_1\wedge\mathfrak f_1,P_1,P_2)\bigr)
\longrightarrow W_1\times W_1
\end{aligned}
\]
以及评价映射
\[
\begin{aligned}
ev_{\mathfrak f_1\wedge\mathfrak f_1,0}\colon
&\bigl(\mathfrak t_{\mathfrak f,2}
\wedge\mathfrak t_{\mathfrak f,2}\bigr)^{-1}(0)\\
&\quad\cap
\bigl(W_1\times W_1\times
U_3(\mathfrak f_1\wedge\mathfrak f_1,P_1,P_2)\bigr)
\longrightarrow L.
\end{aligned}
\]
于是得到
$\mathfrak f_1(P_1)\wedge\mathfrak f_1(P_2)$ 的德拉姆表示
\begin{align}
\mathfrak f_1(P_1)\wedge\mathfrak f_1(P_2)
={}&ev_{\mathfrak f_1!}
\bigl(\pi_{\mathfrak f_1,W_1}^*(\omega_1)\bigr)
\wedge
ev_{\mathfrak f_1!}
\bigl(\pi_{\mathfrak f_1,W_1}^*(\omega_1)\bigr)
\notag\\
={}&ev_{\mathfrak f_1\wedge\mathfrak f_1,0!}
\bigl(
\pi_{\mathfrak f_1\wedge\mathfrak f_1,W_1\times W_1}^*
(\omega_1\wedge\omega_1)
\bigr).
\tag{7.5.28}\label{eq:7-5-28}
\end{align}

定义
\[
\operatorname{comp}\colon
U_2{}_{ev_0}\!\mathbin{\times_{ev_0}}U_2
\longrightarrow L^3
\]
为
\begin{equation}
\operatorname{comp}
\bigl((x_0,x_1),(y_0,y_1)\bigr)
=(x_0,x_1,y_1).
\tag{7.5.29}\label{eq:7-5-29}
\end{equation}
注意 $x_0=y_0$。
可以选择 $U'_3$，使得
\[
U'_3\subset
\operatorname{comp}
\bigl(U_2{}_{ev_0}\!\mathbin{\times_{ev_0}}U_2\bigr).
\]
此时 $\operatorname{comp}$ 是开嵌入，而
\[
U'_3\longrightarrow
\operatorname{comp}
\bigl(U_2{}_{ev_0}\!\mathbin{\times_{ev_0}}U_2\bigr)
\]
也是开嵌入。
所以有开嵌入
\begin{equation}
\varphi\colon U'_3\longrightarrow
U_2{}_{ev_0}\!\mathbin{\times_{ev_0}}U_2.
\tag{7.5.30}\label{eq:7-5-30}
\end{equation}
此外，有覆盖 $\varphi$ 的丛同构
\[
\widetilde\varphi\colon
E_3(\mathfrak f)
\longrightarrow E_2(\mathfrak f)\times E_2(\mathfrak f).
\]
它们诱导交换图
\begin{equation}
\begin{array}{ccc}
E_3(\mathfrak f)
&\xrightarrow{
\ \widetilde\varphi_{\mathfrak f_1\wedge\mathfrak f_1,\mathfrak f_2}\ }
&E_2(\mathfrak f)\times E_2(\mathfrak f)
\\[-0.1em]
\Big\downarrow&&\Big\downarrow
\\[-0.1em]
U_3(\mathfrak f,P_1,P_2)
&\xrightarrow{
\ \varphi_{\mathfrak f_1\wedge\mathfrak f_1,\mathfrak f_2}\ }
&U_3(\mathfrak f_1\wedge\mathfrak f_1,P_1,P_2).
\end{array}
\tag{7.5.31}\label{eq:7-5-31}
\end{equation}
现在利用式~\eqref{eq:7-5-31}，
找出 $\partial_1U_3(\mathfrak f,P_1,P_2)$
的仓西邻域的一个坐标变换。
注意，集合 $\partial_1U_3(\mathfrak f,P_1,P_2)$
在 $U_3(\mathfrak f,P_1,P_2)$ 中的原仓西邻域是
\begin{equation}
(1-\rho,1]\times E_3(\mathfrak f)
\longrightarrow
(1-\rho,1]\times\partial_1U_3(\mathfrak f,P_1,P_2).
\tag{7.5.32.1}\label{eq:7-5-32-1}
\end{equation}
另取仓西坐标图
\begin{equation}
(1-\rho,1]\times E_2(\mathfrak f)\times E_2(\mathfrak f)
\longrightarrow
(1-\rho,1]\times
U_3(\mathfrak f_1\wedge\mathfrak f_1,P_1,P_2).
\tag{7.5.32.2}\label{eq:7-5-32-2}
\end{equation}
于是映射对
\[
\bigl(
\widetilde\varphi_{\mathfrak f_1\wedge\mathfrak f_1,\mathfrak f_2},
\varphi_{\mathfrak f_1\wedge\mathfrak f_1,\mathfrak f_2}
\bigr)
\]
是 $U_3(\mathfrak f,P_1,P_2)$ 中这两个仓西坐标图间的
坐标变换。

现在取一个维数充分大的流形 $W$，
其上带有紧支撑的光滑概率测度，并置
\[
W_2=W_1\times W_1\times W_1\times W.
\]
第一个因子 $W_1$ 将用于
$\partial_0U_3(\mathfrak f,P_1,P_2)$ 附近的扰动，
第二、第三个因子 $W_1\times W_1$ 将用于
$\partial_1U_3(\mathfrak f,P_1,P_2)$ 的一个邻域内。

在仓西坐标图~\eqref{eq:7-5-26-2} 上定义
\[
\begin{aligned}
\mathfrak t_{\mathfrak f_1\circ\mathfrak m_2}^{[0,\rho)'}\colon
&W_2\times[0,\rho)\times
U_3(\mathfrak f_1\circ\mathfrak m_2,P_1,P_2)\\
&\longrightarrow
[0,\rho)\times E_3(\mathfrak m)\times E_2(\mathfrak f)
\end{aligned}
\]
如下：
\begin{align}
&\mathfrak t_{\mathfrak f_1\circ\mathfrak m_2}^{[0,\rho)'}
\bigl((w_{1,1},w_{1,2},w_{1,3},w),a,(\mathbf p,\mathbf q)\bigr)
\notag\\
&\qquad={}
\bigl(a,
\mathfrak s_{\mathfrak m,3}(\mathbf p),
\mathfrak t_{\mathfrak f,2}(w_{1,1},\mathbf q)\bigr).
\tag{7.5.33}\label{eq:7-5-33}
\end{align}
这里 $\mathfrak s_{\mathfrak m,3}$ 是定义 $\mathfrak m_2$
所用的截面，$\mathfrak t_{\mathfrak f,2}$ 是定义
$\mathfrak f_1$ 所用的 $W_1$ 参数截面族；
并取
$\mathbf p\in U_3(\mathfrak m,P_1,P_2)$、
$\mathbf q\in U_2(\mathfrak f)$
（见式~\eqref{eq:7-5-19}）。

接着在仓西坐标图~\eqref{eq:7-5-32-2} 上定义
\[
\begin{aligned}
\mathfrak t_{\mathfrak f_1\wedge\mathfrak f_1}^{(1-\rho,1]'}\colon
&W_2\times(1-\rho,1]\times
U_3(\mathfrak f_1\wedge\mathfrak f_1,P_1,P_2)\\
&\longrightarrow
(1-\rho,1]\times E_2(\mathfrak f)\times E_2(\mathfrak f)
\end{aligned}
\]
为
\begin{align}
&\mathfrak t_{\mathfrak f_1\wedge\mathfrak f_1}^{(1-\rho,1]'}
\bigl((w_{1,1},w_{1,2},w_{1,3},w),a,
(\mathbf p_1,\mathbf p_2)\bigr)
\notag\\
&\qquad={}
\bigl(a,
\mathfrak t_{\mathfrak f,2}(w_{1,2},\mathbf p_1),
\mathfrak t_{\mathfrak f,2}(w_{1,3},\mathbf p_2)\bigr).
\tag{7.5.34}\label{eq:7-5-34}
\end{align}
为了把式~\eqref{eq:7-5-33} 与~\eqref{eq:7-5-34}
延拓到 $W_2\times U_3(\mathfrak f,P_1,P_2)$ 的仓西邻域，
需要修改它们，使它们同式~\eqref{eq:7-5-26-1}、
\eqref{eq:7-5-32-1} 中将给出的某些截面相容。
具体地，需要下面的引理。

\begin{lemma}{7.5.35}\label{lem:7-5-35}
可以把
$\mathfrak t_{\mathfrak f_1\circ\mathfrak m_2}^{[0,\rho)'}$
修改为
$\mathfrak t_{\mathfrak f_1\circ\mathfrak m_2}^{[0,\rho)}$，
使得：
\begin{align}
\mathfrak t_{\mathfrak f_1\circ\mathfrak m_2}^{[0,\rho)}
&=
\mathfrak t_{\mathfrak f_1\circ\mathfrak m_2}^{[0,\rho)'}
\quad\text{在 }W_2\times[0,\rho/3)\times
U_3(\mathfrak f_1\circ\mathfrak m_2,P_1,P_2)\text{ 上};
\tag{7.5.36.1}\label{eq:7-5-36-1}
\end{align}
\textup{若}
\begin{equation*}
\mathfrak t_{\mathfrak f_1\circ\mathfrak m_2}^{[0,\rho)}
(\vec w,a,\mathbf p)=0,
\qquad a\in(2\rho/3,\rho),
\end{equation*}
\textup{则存在}
\begin{equation}
\mathbf p'\in\partial_0U_3(\mathfrak f,P_1,P_2),
\qquad
\varphi_{\mathfrak f_1\circ\mathfrak m_2,\mathfrak f_2}
(\mathbf p')=\mathbf p.
\tag{7.5.36.2}\label{eq:7-5-36-2}
\end{equation}
\textup{存在截面}
\begin{equation}
\begin{aligned}
\mathfrak t_{\mathfrak f,2}^{(2\rho/3,\rho)}\colon
&W_2\times(2\rho/3,\rho)\times
\partial_0U_3(\mathfrak f,P_1,P_2)\\
&\longrightarrow(2\rho/3,\rho)\times E_3(\mathfrak f),
\end{aligned}
\tag{7.5.36.3}\label{eq:7-5-36-3}
\end{equation}
\textup{使得}
$\mathfrak t_{\mathfrak f,2}^{(2\rho/3,\rho)}(\vec w,\cdot)$
\textup{关于坐标变换}
\[
\bigl(
\widetilde\varphi_{\mathfrak f_1\circ\mathfrak m_2,\mathfrak f_2},
\varphi_{\mathfrak f_1\circ\mathfrak m_2,\mathfrak f_2}
\bigr)
\]
\textup{同}
$\mathfrak t_{\mathfrak f_1\circ\mathfrak m_2}^{[0,\rho)}
(\vec w,\cdot)$
\textup{在定义 A1.21 的意义下相容。}

\textup{最后，}
$\mathfrak t_{\mathfrak f_1\circ\mathfrak m_2}^{[0,\rho)}$
\textup{具有如下横截性：集合}
\begin{equation}
\begin{aligned}
\mathfrak X_{[0,\rho)}
=\bigl\{(\vec w,a,\mathbf p)&\in
W_2\times[0,\rho)\times
U_3(\mathfrak f_1\circ\mathfrak m_2,P_1,P_2)\ \big|\\
&\mathfrak t_{\mathfrak f_1\circ\mathfrak m_2}^{[0,\rho)}
(\vec w,a,\mathbf p)=0\bigr\}
\end{aligned}
\tag{7.5.36.4}\label{eq:7-5-36-4}
\end{equation}
\textup{是子流形，且评价映射}
$ev_0\colon\mathfrak X_{[0,\rho)}\to L$
\textup{是浸没。这里的 $ev_0$ 由}
\[
ev_0\colon
U_3(\mathfrak f_1\circ\mathfrak m_2,P_1,P_2)
\longrightarrow L
\]
\textup{诱导}\translatornote{源书 (7.5.36.2) 把
$\partial_0U_3(\mathfrak f,P_1,P_2)$ 误写成
$\partial_0U_3(\mathfrak f,P_2,P_2)$；
(7.5.36.4) 的集合元素先写 $(\vec w,a,x)$，
零点条件却使用 $\mathbf p$。
译文分别恢复 $P_1$ 与统一变量 $\mathbf p$。
见疑误 A-0622。}。
\end{lemma}

\begin{proof}
先用标准横截性论证构造扰动
$\mathfrak t_{\mathfrak f,2}^{(2\rho/3,\rho)}$。
再要求式~\eqref{eq:7-5-36-2}、
\eqref{eq:7-5-36-3} 成立，
从而在
\[
W_2\times(2\rho/3,\rho)\times
\partial_0U_3(\mathfrak f,P_1,P_2)
\]
上定义
$\mathfrak t_{\mathfrak f_1\circ\mathfrak m_2}^{[0,\rho)}$。
在
$W_2\times(0,\rho/3)\times
U_3(\mathfrak f_1\circ\mathfrak m_2,P_1,P_2)$ 上，
该截面由式~\eqref{eq:7-5-36-1} 决定。
现在再用标准横截性论证，把它们延拓到
\[
W_2\times[0,\rho)\times
\partial_0U_3(\mathfrak f,P_1,P_2).
\]
\end{proof}

同样可以把
$\mathfrak t_{\mathfrak f_1\wedge\mathfrak f_1}^{(1-\rho,1]'}$
修改为
$\mathfrak t_{\mathfrak f_1\wedge\mathfrak f_1}^{(1-\rho,1]}$，
使截面
\[
\begin{aligned}
\mathfrak t_{\mathfrak f,2}^{(1-\rho,1-2\rho/3)}\colon
&W_2\times(1-\rho,1-2\rho/3)\times
\partial_1U_3(\mathfrak f,P_1,P_2)\\
&\longrightarrow
(1-\rho,1-2\rho/3)\times E_3(\mathfrak f)
\end{aligned}
\]
同
$\mathfrak t_{\mathfrak f_1\wedge\mathfrak f_1}^{(1-\rho,1]}$
在定义 A1.21 的意义下相容。
还要求
\[
\mathfrak t_{\mathfrak f_1\wedge\mathfrak f_1}^{(1-\rho,1]'}
=\mathfrak t_{\mathfrak f_1\wedge\mathfrak f_1}^{(1-\rho,1]}
\]
在
\[
W_2\times(1-\rho/2,1]\times
U_3(\mathfrak f_1\wedge\mathfrak f_1,P_1,P_2)
\]
上成立，并要求
\[
\begin{aligned}
\mathfrak X_{(1-\rho,1]}
=\bigl\{(\vec w,a,\mathbf p)&\in
W_2\times(1-\rho,1]\times
U_3(\mathfrak f_1\wedge\mathfrak f_1,P_1,P_2)\ \big|\\
&\mathfrak t_{\mathfrak f_1\wedge\mathfrak f_1}^{(1-\rho,1]}
(\vec w,a,\mathbf p)=0\bigr\}
\end{aligned}
\]
为子流形，且评价映射
$ev_0\colon\mathfrak X_{(1-\rho,1]}\to L$ 为浸没
\translatornote{源书把上述相等式的定义域写成
$U_2(\mathfrak f_1\wedge\mathfrak f_1;P_1,P_2)$，
但该空间在本节定义并始终记作 $U_3$；
又在 $\mathfrak X_{(1-\rho,1]}$ 的零点条件中误用了
$\mathfrak t_{\mathfrak f_1\circ\mathfrak m_2}^{[0,\rho)}$。
译文按当前 $\partial_1$ 端的仓西坐标图及扰动改正。
见疑误 A-0623。}。

注意
\[
U_3(\mathfrak f,P_1,P_2)
\cong[0,1]\times\partial_0U_3(\mathfrak f,P_1,P_2)
\cong[0,1]\times\partial_1U_3(\mathfrak f,P_1,P_2).
\]
因此，
$\mathfrak t_{\mathfrak f,2}^{(2\rho/3,\rho)}$ 与
$\mathfrak t_{\mathfrak f,2}^{(1-\rho,1-2\rho/3)}$
在
\[
W_2\times
\bigl((2\rho/3,\rho)\cup(1-\rho,1-2\rho/3)\bigr)
\times\partial_0U_3(\mathfrak f,P_1,P_2)
\]
上定义一个截面。
在下一个引理中，把
\[
(2\rho/3,1-2\rho/3)
\times\partial_0U_3(\mathfrak f,P_1,P_2)
\]
看作 $U_3(\mathfrak f,P_1,P_2)$ 的子集。

\begin{lemma}{7.5.37}\label{lem:7-5-37}
存在截面
\[
\begin{aligned}
\mathfrak t_{\mathfrak f,2}^{(2\rho/3,1-2\rho/3)}\colon
&W_2\times(2\rho/3,1-2\rho/3)
\times\partial_0U_3(\mathfrak f,P_1,P_2)\\
&\longrightarrow
(2\rho/3,1-2\rho/3)\times E_3(\mathfrak f),
\end{aligned}
\]
并具有下列性质：

\textup{它在}
\begin{equation}
W_2\times
\bigl((3\rho/4,\rho)\cup(1-\rho,1-3\rho/4)\bigr)
\times\partial_0U_3(\mathfrak f,P_1,P_2)
\tag{7.5.38.1}\label{eq:7-5-38-1}
\end{equation}
\textup{上分别同}
$\mathfrak t_{\mathfrak f,2}^{(2\rho/3,\rho)}$、
$\mathfrak t_{\mathfrak f,2}^{(1-\rho,1-2\rho/3)}$
\textup{一致。}

$\mathfrak t_{\mathfrak f,2}^{(2\rho/3,1-2\rho/3)}$
\textup{具有如下横截性：集合}
\begin{equation}
\begin{aligned}
\mathfrak X_{(2\rho/3,1-2\rho/3)}
=\bigl\{(\vec w,a,\mathbf p)&\in
W_2\times(2\rho/3,1-2\rho/3)
\times\partial_0U_3(\mathfrak f,P_1,P_2)\ \big|\\
&\mathfrak t_{\mathfrak f,2}^{(2\rho/3,1-2\rho/3)}
(\vec w,a,\mathbf p)=0\bigr\}
\end{aligned}
\tag{7.5.38.2}\label{eq:7-5-38-2}
\end{equation}
\textup{是子流形，且评价映射}
\[
ev_0\colon\mathfrak X_{(2\rho/3,1-2\rho/3)}\longrightarrow L
\]
\textup{是浸没。}
\translatornote{源书 (7.5.38.2) 的零点条件把截面下标
$\mathfrak f,2$ 漏作单独的 $2$；译文依引理陈述及
同一式左端补回。见疑误 A-0624。}
\end{lemma}

证明仍是横截性定理的一个标准应用。
由式~\eqref{eq:7-5-38-1}、\eqref{eq:7-5-36-3}，
可以把
\[
\left.
\mathfrak t_{\mathfrak f,2}^{(2\rho/3,1-2\rho/3)}
\right|_{
W_2\times((3\rho/4,\rho)\cup(1-\rho,1-3\rho/4))
\times\partial_0U_3(\mathfrak f,P_1,P_2)}
\]
同
$\mathfrak t_{\mathfrak f_1\circ\mathfrak m_2}^{[0,\rho)}$
胶合。
对
$\mathfrak t_{\mathfrak f_1\wedge\mathfrak f_1}^{(1-\rho,1]}$
也有类似说法。
定义
\[
\begin{aligned}
\mathfrak X
={}&\mathfrak X_{[0,\rho)}\cup\mathfrak X_{(1-\rho,1]}\\
&\cup\Bigl(
\mathfrak X_{(2\rho/3,1-2\rho/3)}
\cap\bigl(
W_2\times(3\rho/4,1-3\rho/4)
\times\partial_0U_3(\mathfrak f,P_1,P_2)
\bigr)\Bigr).
\end{aligned}
\]
按定义，$ev_0\colon\mathfrak X\to L$ 是浸没；
另有真投影
\[
\pi_{W_2}\colon\mathfrak X\longrightarrow W_2.
\]
现在可以定义 $\mathfrak f_2(P_1,P_2)$。

\begin{definition}{7.5.39}\label{def:7-5-39}
用如下德拉姆表示定义
\translatornote{源书刚把评价映射定义为
$ev_0\colon\mathfrak X\to L$，定义 7.5.39 的公式却写
$ev_!$；译文补回下标 $0$。见疑误 A-0625。}
$\mathfrak f_2(P_1,P_2)$：
\[
\mathfrak f_2(P_1,P_2)
=ev_{0!}\bigl(\pi_{W_2}^*\omega_2\bigr),
\]
其中 $\omega_2$ 是固定在 $W_2$ 上的光滑概率测度，
它是 $W_1$ 与 $W$ 上诸测度的直积。
\end{definition}

\begin{proposition}{7.5.40}\label{prop:7-5-40}
按上述方式定义 $\mathfrak f_2(P_1,P_2)$ 后，
式~\eqref{eq:7-5-16} 成立。
\end{proposition}

\begin{proof}
这是式~\eqref{eq:7-5-17}、\eqref{eq:7-5-20}、
\eqref{eq:7-5-28}、\eqref{eq:7-5-33}、
\eqref{eq:7-5-36-1} 及引理~\ref{lem:7-5-10} 的推论。
\end{proof}
